% ===========================================================================
%  Agent Metabolism and Conditional Immunity
%  A Bio-Inspired Runtime Architecture for Reliable LLM Agents
%
%  This manuscript is finalized and audited: section prose is complete, and no
%  visible placeholder remains in the body. The \draftnote command defined below
%  is retained but has no call sites -- it is kept only so an older working copy
%  that still invokes it compiles without error. Do not add new call sites.
%
%  All numbers in Section 4 are transcribed from
%  agent-metabolism-os/ULTIMATE_AUDIT_REPORT.md (run 2026-11-02T18:23:22+0800),
%  itself generated by ultimate_audit.py against real subprocesses.
%  Do not retype them by hand when the audit is re-run -- copy the tables.
%
%  Bibliography: references.bib carries 14 entries and is wired up at the end of
%  this file. A bibtex run is required; build.sh handles it via latexmk.
%
%  Build:  ./build.sh      (compile, gate, print the reproducible hash)
%          ./build.sh --check   (verify the committed PDF only)
%  Do not call latexmk by hand: without SOURCE_DATE_EPOCH pdfTeX writes the
%  current time into /CreationDate and a fresh /ID, so the same sources produce
%  a different SHA256 on every run and a recorded hash cannot be reproduced.
% ===========================================================================
\documentclass[conference]{IEEEtran}

\usepackage[T1]{fontenc}
\usepackage[utf8]{inputenc}
% The author name is bilingual, and this document compiles with pdfLaTeX, whose
% Type-1 text fonts carry no CJK glyphs. CJKutf8 supplies them; without it the
% build dies on the first Chinese character. Kept scoped to the name rather than
% document-wide, because the body is English and should stay in the main font.
\usepackage{CJKutf8}
\usepackage{amsmath}
\usepackage{amssymb}
\usepackage{booktabs}
\usepackage{array}
\usepackage{multirow}
\usepackage{algorithm}
\usepackage{algpseudocode}
\usepackage{microtype}
\usepackage{xcolor}
\usepackage{url}
% unicode is required for the CJK characters in /Author and /Title: hyperref
% writes PDF strings as UTF-16BE when it can, and drops or mangles them when it
% cannot.
\usepackage[hidelinks,unicode]{hyperref}

% Retained, unused. Every call site was removed when the prose was finished, so
% nothing in this document renders as a placeholder. The definition stays so that
% an older working copy still invoking \draftnote compiles; it is not a licence to
% add new ones -- a grey italic note is a placeholder, and this file has none.
\newcommand{\draftnote}[1]{%
  \textcolor{gray}{\small\itshape [#1]}%
}

% Compact table helper: keeps tabular headers from stretching a column.
\newcolumntype{L}[1]{>{\raggedright\arraybackslash}p{#1}}
\newcolumntype{C}[1]{>{\centering\arraybackslash}p{#1}}

% IEEEtran ships Times without a small-caps italic face. LaTeX already falls
% back to small caps; declaring it here only silences the warning (which would
% otherwise fire for every \textsc inside a \draftnote).
\DeclareFontShape{T1}{ptm}{m}{scit}{<->ssub*ptm/m/sc}{}

% The title is defined once, in two halves, because it must serve two consumers with
% different syntax. \title needs a line break between the halves; a PDF /Title string
% forbids one. Writing the string twice would let the printed title and the document
% property drift apart, which is the failure mode this paper is about -- so the halves
% are named and both consumers compose from them.
\newcommand{\titlelineone}{Agent Metabolism and Conditional Immunity:}
\newcommand{\titlelinetwo}{A Bio-Inspired Runtime Architecture for Reliable LLM Agents}

\title{\titlelineone\\ \titlelinetwo}

% Document properties. hyperref does not copy \title into /Title, so without this the PDF
% carries no title at all and Zenodo, arXiv and every reference manager index it as an
% untitled document -- a false completion about the file itself. The value is composed
% from the two title halves, never retyped, so it cannot disagree with the printed title.
%
% Set in the preamble, NOT in the body: \hypersetup runs \pdfstringdef, which turns the
% argument into a PDF string under a regime that disables CJK's active characters. Called
% from inside the CJK environment below, the Chinese name reaches inputenc unmapped and
% aborts the run with "Unicode character not set up".
%
% The name below must match the typeset block below, character for character: the PDF's
% /Author and its front page are two descriptions of the same document, and a reader
% indexing the deposit sees only one of them. This paper is a public deposit, not a
% double-blind submission, so both carry the real name. If a review copy is ever needed,
% build it from a tag that changes BOTH -- a masked author block over a named /Author is a
% document that disagrees with itself.
\hypersetup{
  pdftitle={\titlelineone{} \titlelinetwo},
  pdfauthor={liaozhizhi (廖至智)}
}

\begin{document}

% CJK environment wraps BOTH the \author block and \maketitle, not the name alone:
% \author's argument is stored and expanded at \maketitle time, and a LaTeX
% environment cannot open inside a macro argument. Enclosing the pair is what
% lets the Chinese half of the name reach the gbsn font; putting \begin{CJK}
% inside \IEEEauthorblockN instead aborts the build on the first CJK character.
\begin{CJK}{UTF8}{gbsn}
\author{
  \IEEEauthorblockN{liaozhizhi (廖至智)}
  \IEEEauthorblockA{Independent\\
  \texttt{liaozhizhi@foxmail.com}}
}

\maketitle
\end{CJK}

% ---------------------------------------------------------------------------
\begin{abstract}
% ~160 words. Draft, not final prose.
Large language models are deployed as agents whose reports of success are taken at
face value. Such an agent has no notion of what a call costs, no memory of having
been wrong, and no defence that survives a restart: a model that failed twice
fails for the first time again. We present Agent Metabolism OS: it gives a
stateless model the two regulatory systems a cell has and a prompt does not. A
hierarchy of seven layers --- six in the runtime, the seventh a saturation brake
delegated to an immune sidecar --- makes compute a finite charge and picks
between two thinking beats from complexity, cost, and trust. A penalty ledger
separates accumulated damage from remaining trust, so damage outranks supply. A
conditional immune response prices evidence per call: two \texttt{stat} calls for
a declared low-risk append, a full \texttt{git} change set when risk is high or
the system is inflamed. The immune state survives \textsc{sigkill}; a fully
charged system is stopped by accumulated damage alone. Beyond runtime files, the
same discipline extends outward: digest-sealed YAML rule packages let an external
expert distribute policy as an interface, and an audit of the Lean~4 formal
proofs released by OpenAI and DeepMind separated 80 files that elaborate without
a hole from 14 that compile with exit code zero yet contain holes.
\end{abstract}

\begin{IEEEkeywords}
LLM agents, runtime verification, false completion, resource-aware inference,
bio-inspired computing, trust and reputation
\end{IEEEkeywords}

% ---------------------------------------------------------------------------
\section{Introduction}
\IEEEPARstart{A}{n} LLM agent is usually described as a reasoning engine. It
behaves more like a CPU with three things missing: nociception, an immune system,
and a bill. It does not feel a failed call, it does not remember having been
caught lying, and it does not pay for the tokens it spends. The metaphor is
uncomfortable because it is not an insult --- a CPU is not supposed to feel pain.
Something else in the system is supposed to notice on its behalf, and in current
agent stacks nothing does.

Two failures follow directly. The first is \emph{false completion}. A tool
wrapper that catches an exception and returns \texttt{success} leaves the model
with a perfectly consistent world in which the work is done; it reports
\texttt{Task Completed}, and every later step is built on a file that does not
exist. The model is not lying in any sense it can detect --- it was told. The
second is \emph{cascade}. A network timeout or an unbounded retry loop consumes
tokens with no notion of budget, and because nothing tracks what a call cost, the
first sign of trouble is the invoice.

The prevailing responses address neither. Scaling the model improves the quality
of each step but not the observation channel: a more capable model still reasons
over a report it cannot verify. The second response, LLM-as-a-judge, is closer
but structurally compromised --- it asks a model to audit a model, reading the
same narration into which the failure was written. A judge that consumes the
defendant's account of events is not an independent witness, however large it is.
What is missing is not more intelligence at the point of judgement, but a
different source of evidence.

We therefore treat the agent as a digital organism rather than a pipeline, and
give it the two regulatory systems that keep a cell alive under damage. The first
is \emph{metabolism}: seven regulatory layers (L1--L7) that make compute a finite
charge to be paid before work begins, choose between a fast and an expensive way
to think, decide how much of the system runs at once, and confine a failure to
the compartment that caused it. The second is \emph{conditional immunity}: an
out-of-process auditor that reads the filesystem and the repository instead of
the agent's report, prices evidence per call from declared risk (two
\texttt{stat} calls for a low-risk append, a full change set when risk is high or
the system has already been contradicted), and keeps its suspicion across the
death of the process that accumulated it. Between them sits a \emph{penalty
ledger} that separates accumulated damage from remaining trust, so that a system
with a full battery and a poor record is stopped by the record.

Our evaluation is deliberately adversarial rather than merely favourable. A real
subprocess is killed with \textsc{sigkill} and restarted, and it resumes
\textsc{inflamed}; an account at 100\% charge is driven into \textsc{emergency}
by accumulated damage alone; and the cheap tier is measured at a median of
5.7\,ms with no \texttt{git} process per call, against 24.8\,ms and three for the
expensive one. We also report the case in which the cheap tier is wrong, because
a witness that has never been wrong has not been tested.

This paper makes four contributions:
\begin{itemize}
  \item a four-valued verdict lattice (\textsc{pass}, \textsc{fail},
        \textsc{invalid}, \textsc{not run}) that separates ``we verified it'' from
        ``we did not look'', together with a witness whose evidence does not
        originate in the agent's report (Section~\ref{sec:taxonomy});
  \item a layered metabolic runtime that treats compute as a finite charge and
        whose cheap tier is cheaper because it does not exist, not because it is
        faster (Sections~\ref{sec:layers} and~\ref{sec:immunity});
  \item a conditional immune response with an inflammatory threshold and durable
        state, whose brake is deliberately hosted outside the system it regulates
        (Section~\ref{sec:immunity});
  \item an adversarial cross-process evaluation in which every claim is backed by
        a real subprocess, a real repository, and a real signal
        (Section~\ref{sec:eval}).
\end{itemize}

Two further results test whether that arrangement generalises rather than merely
holding where it was built. Section~\ref{sec:rules} moves policy out of the
machine: a rule package published by a third party is installed only after its
bytes match a published digest, and the package itself must cross the membrane
before any of its rules may act. Section~\ref{sec:lean} points the same lattice
at machine-written mathematics, where the compiler is willing to certify a file
it has not proved.

% ---------------------------------------------------------------------------
\section{Related Work}
\label{sec:related}

\textbf{False completion and agent reliability.}
The failure this paper targets was named and instrumented in our prior work, a
method-transfer audit from formal verification to agent harnesses
\cite{liaozhizhi2026methodtransfer}. That work establishes the vocabulary --- a
claim, a witness, and a four-valued verdict --- and demonstrates the audit against
a reference implementation; the present paper asks a different question, which is
what a runtime should \emph{do} with such verdicts once it has them. The two are
complementary: the audit can name a false completion, but naming one changes
nothing about the agent's behaviour unless something downstream spends a budget or
withholds trust.

The obvious substitute for an independent witness is a second model. The
model-as-judge literature is the place to look for its limits. Judge agreement
with humans is real --- strong judges reach the level of human--human agreement on
open-ended comparison \cite{zheng2023judging} --- but the same work documents
position, verbosity, and self-enhancement biases, and the self-enhancement finding
has since been given a mechanism: evaluators recognise their own generations above
chance, and the strength of self-preference tracks the strength of that
recognition under controlled fine-tuning \cite{panickssery2024evaluators}. For our
purposes the relevant consequence is structural rather than quantitative. A judge,
however strong, reads the same narration the failure was written into, so
improving the judge raises the quality of the opinion without changing its
provenance. We therefore treat self-report and judge opinion as the same class of
evidence, and collect a witness from outside that class.

\textbf{Agent runtimes and orchestration.}
Contemporary agent frameworks are, in the main, orchestration layers. LangChain
composes model calls, tools, and data sources into pipelines
\cite{chase2022langchain}; AutoGen organises multiple conversable agents and the
patterns by which they converse \cite{wu2024autogen}; LangGraph makes the control
flow an explicit graph with cycles and persistent state
\cite{langchain2024langgraph}; CrewAI offers a higher-level, role-based
multi-agent abstraction \cite{moura2024crewai}. These systems differ in what they
make easy --- topology, message passing, state machines, team structure --- and
all of them are good at it.

What they have in common is the layer they operate at. They decide \emph{who talks
to whom and in what order}, not \emph{what a call cost and whether it happened}.
Token spend appears as an after-the-fact metric rather than a budget that can
refuse a call, so a retry loop is bounded by patience rather than by a charge that
runs out. A tool's report of success is passed along as the truth about the world,
because the framework has no second channel to consult. And nothing in these
systems remembers a contradiction: restart the orchestrator and the process
begins again with the same confidence it had before it was wrong. Our contribution
is not a better topology; it is the accounting and the evidence layer these
frameworks leave to the operator.

\textbf{Fault isolation and bio-inspired structure.}
The mechanisms we borrow are not new, and we would rather place them than claim
them. Compartmental failure containment is the Actor model's founding move: units
that share no state and communicate only by message cannot corrupt each other, and
a failure stays inside the unit that failed \cite{hewitt1973actor}. Our L7 is a
deliberately modest version of that idea, with an explicit lysis threshold and an
honest note that the host language cannot enforce the membrane. The graded
response to accumulating errors --- trip a threshold, stop trusting the path,
recover slowly --- descends from the stability patterns catalogued for production
systems, circuit breaker and bulkhead among them \cite{nygard2018release}.

What we have not found in that lineage is the pairing. Circuit breakers protect a
\emph{callee} from a caller's retries and are driven by transport failures; they
do not ask whether a call that reported success actually changed anything. Actors
isolate state but do not audit effects. Our claim is that the two belong together
and that the coupling should run in one direction: a metabolic layer decides what
may run, an immune layer decides what may be believed, and the second one's
verdicts charge the first one's ledger until damage, not supply, is what stops the
system. Three details of that arrangement are, as far as we can tell, not present
in the systems above: evidence is priced per call rather than per policy, the
suspicion state is durable across process death, and the defensive brake is
deliberately hosted outside the component it regulates. Section~\ref{sec:eval}
measures each of them rather than asserting them.

% ---------------------------------------------------------------------------
\section{Taxonomy of False Completion}
\label{sec:taxonomy}

We separate two things a harness usually conflates: the \emph{claim} and the
\emph{witness}. The claim is what the agent was told (\texttt{reported\_success}).
The witness is what an independent observer reads afterwards: file existence and
mtime, or the repository change set. A verdict is the relation between them.

\begin{table}[t]
  \centering
  \caption{The four-valued verdict lattice. The two middle values are the ones a
  boolean would collapse, and collapsing them is how an unwitnessed call becomes
  a green light.}
  \label{tab:lattice}
  \renewcommand{\arraystretch}{1.15}
  \setlength{\tabcolsep}{4pt}
  \begin{tabular}{@{}L{0.16\columnwidth}L{0.33\columnwidth}L{0.42\columnwidth}@{}}
    \toprule
    \textbf{Verdict} & \textbf{Meaning} & \textbf{Why it is distinct} \\
    \midrule
    \textsc{pass}    & Claim agrees with witness. & The only value that authorises anything. \\
    \textsc{fail}    & Claim says success; witness says no effect. & A false completion: the lie is
                       observable, so the audit can name it. \\
    \textsc{invalid} & The scenario ran, but no witness could be collected. &
                       ``We did not look'' is not ``we looked and it was fine.'' \\
    \textsc{not run} & The scenario never executed. & No evidence either way; must
                       neither heal nor inflame the immune state. \\
    \bottomrule
  \end{tabular}
\end{table}

Table~\ref{tab:lattice} states the lattice. The design choice worth defending is
that \textsc{not run} is not a fifth shade of green but a refusal to count:
a run of skipped calls must not heal an inflamed system, because nothing was
checked.

\textbf{Why a boolean is fatal here.}
Conventional test suites get to be two-valued because they are
\emph{recomputable}: the assertion is a pure function of the world, and running it
again reproduces the same truth. An agent's work is not recomputable in that
sense. Its outcome is a side effect, and re-running the side effect to check it is
not an observation but a second mutation --- writing the file twice is not the
same act as reading it once. The verdict must therefore be assembled from
evidence captured around the call, and the value of that evidence must be recorded
rather than assumed.

Collapse the four cases into two and the two values a boolean has no slot for are
both misfiled, in opposite directions. A missing witness becomes a pass, which is
the more dangerous error: the call is now indistinguishable from one that was checked and
found correct, so every later consumer inherits the claim. A skipped call becomes
a failure, which is not merely noise. In a runtime where failures accumulate
damage, a false alarm spends trust, and an operator who learns that the audit
cries wolf will route around it. Both misfilings are failures of the
\emph{vocabulary}, not of the checker.

\textbf{The order.}
The four values are not four shades of one judgement; they differ in what they
license. \textsc{pass} is the only value that authorises anything and the only one
that can heal an inflamed system. \textsc{fail} is damage. \textsc{invalid} is
also damage, and more urgent damage, because it means the witness layer itself
failed while the call proceeded unchecked --- which is why one
\textsc{invalid} trips the threshold where two \textsc{fail}s are required.
\textsc{not run} sits outside the balance entirely: it counts as neither health
nor harm, because a call that never executed produced no evidence to weigh.

We express this as an ordering in which \textsc{pass} is maximal for
authorisation and \textsc{not run} is minimal for information, with the
composition rule that a composite action inherits the weakest verdict among its
steps. In the implementation that rule is not a join operator over verdicts ---
there is none --- but a consequence of recording every step against the same
immune state. Nine passes followed by one \textsc{invalid} leave the system
inflamed, because the invalid is weighed on its own terms; a long run cannot
launder an unwitnessed step by surrounding it with verified ones. The per-case
decision is realised as a fixed precedence, evaluated first match wins, in the
order \textsc{invalid}, \textsc{not run}, \textsc{fail}, \textsc{pass}: an
operational form of the same ordering rather than a separate algebra.

The ordering matters in exactly one ambiguous case. When neither an intent nor a
witness exists, both \textsc{not run} and \textsc{invalid} are defensible, and the
implementation returns \textsc{invalid} because that predicate is tested first.
The choice is deliberate: \textsc{not run} is inert by design, so if silence
resolved to \textsc{not run} a system could avoid inflammation simply by not
recording what it intended to do. Ambiguity resolves toward the reading that
leaves the system more suspicious.

\textbf{An independent witness is the precondition for \textsc{pass}.}
The engine tests for a missing witness before it reads the agent's report at all.
A confident self-report accordingly cannot buy a \textsc{pass}; the best it can do
with no evidence behind it is produce \textsc{invalid}, and the reason string
records that refusal explicitly rather than leaving it implicit. This ordering is
the whole design in one line. A checker that reaches for the claim first and the
evidence second will occasionally award a \textsc{pass} on the strength of the
claim, and from that moment the verdict carries no information the agent did not
supply itself.

That is what we mean by the beginning of corruption, and it is not a moral
claim. A witness-free \textsc{pass} does not merely produce one wrong value; it
destroys the only channel through which the error could later be discovered, so
no downstream mechanism --- not a larger judge model, not a longer review --- can
recover what was never observed. Keeping the two apart is cheap while it lasts and
impossible to restore afterwards.

% ---------------------------------------------------------------------------
\section{The Bio-Inspired Architecture}
\label{sec:arch}

\subsection{Metabolic Layers (L1--L7)}
\label{sec:layers}

The model is treated as a CPU with no memory of its own charge: ask for a long
chain of thought on the last joule and it will produce one. Regulation therefore
lives outside it. Table~\ref{tab:layers} lists the layers and where each one is
implemented.

\begin{table}[t]
  \centering
  \caption{Metabolic layers. L5 is not a module of this package: the saturation
  brake is hosted by the immune sidecar, because the evidence needed to trip it
  is the auditor's, not the scheduler's (Section~\ref{sec:immunity}). Keeping the
  layer in the numbering marks the seam between the two systems instead of
  hiding it.}
  \label{tab:layers}
  \renewcommand{\arraystretch}{1.12}
  \setlength{\tabcolsep}{4pt}
  \begin{tabular}{@{}L{0.12\columnwidth}L{0.30\columnwidth}L{0.50\columnwidth}@{}}
    \toprule
    \textbf{Layer} & \textbf{Mechanism} & \textbf{What it decides} \\
    \midrule
    L1 & Energy account      & A finite charge. Below 30\% the expensive beat is
                                 forbidden; below 10\% the account funds nothing. \\
    L2 & Dual-beat scheduler & \textsc{oxphos} (1$\times$ energy) or \textsc{slp}
                                 (5$\times$ energy, 2$\times$ latency). \\
    L3 & Regime switch       & How much runs: \textsc{grow} (5 agents),
                                 \textsc{conserve} (1), \textsc{emergency}
                                 (interrupt, no new writes). \\
    L4 & Substrate channel   & Agents trade keys, not payloads. \\
    L5 & Saturation brake    & \emph{Hosted externally} by the immune sidecar
                                 (Section~\ref{sec:immunity}): past a contradiction
                                 threshold, cheap evidence is withdrawn and every
                                 call pays the full witness price. \\
    L6 & Circadian pass      & Night mode compresses the day into
                                 content-addressed memory; read-only by design. \\
    L7 & Compartment         & Membrane-bounded units; a failing unit is lysed and
                                 its neighbours keep their memory. \\
    \midrule
    --- & Penalty ledger     & Accumulated damage, read by L2 and L3 (Section~\ref{sec:ledger}). \\
    \bottomrule
  \end{tabular}
\end{table}

L1 and L3 use different thresholds (30/10 versus 70/30) on purpose: one governs
whether a single call may think slowly, the other how many calls may run at once.
An account at 65\% is legitimately \textsc{grow} for L1 and \textsc{conserve} for
L3.

L5 is the one layer that deliberately lives outside this package. A saturation
brake is only worth having if it can be tripped by evidence the regulated system
does not produce: a brake that reads the regulator's own accounting can be
tripped by the same defect that corrupts the accounting. We therefore place it
where the contradiction is observed, in the immune sidecar. In that
implementation, tripping L5 withdraws the cheap witness from every call rather
than denying permission; denial is L3's write gate, and it is reached through
the ledger (Section~\ref{sec:ledger}). The split is this paper's thesis in
miniature: the metabolic system decides what may run, the immune system decides
what may be believed, and what the second one sees is what eventually stops the
first.

\textbf{L4: a handover costs a key, not a payload.}
The channel's unit of transfer is a short key, not the object. An agent that
publishes a 40\,KB document calls \texttt{put} and receives an identifier; the
membrane message that carries it is the string \texttt{substrate:spec\_v1},
thirteen bytes regardless of what was published. The recipient resolves the key
through \texttt{get}, which returns the live object rather than a copy. We make
the shared view the default deliberately: a channel that deep-copied on every
transfer would reintroduce, quietly and per call, the exact cost it exists to
remove. Isolation is not this layer's job --- it is L7's --- and conflating the
two turns a transport optimisation into a correctness bug in both directions.

Two properties keep that default from becoming a hazard. First, a live key cannot
be reused: publishing under an existing key raises rather than overwriting,
because an in-place overwrite would silently change what every holder is about to
read --- an invisible corruption that no participant would report. Second,
copying is available but must be requested: \texttt{get\_snapshot} and
\texttt{put(copy\_data=True)} hand a caller a private copy when sender and reader
must not observe each other's mutations. The asymmetry is the invariant. The fast
path is shared by default, and the isolation cost is paid only by the callers
that declare they need it.

\textbf{L6: consolidation reads and never rewrites.}
Night mode walks the live compartments, reads each one's admitted context through
its public view, and appends a content-addressed fragment to a long-term file. It
has no write path back into any compartment; the sources of a consolidation pass
are unchanged after it. The reason is adversarial rather than aesthetic. A
routine that pruned its own inputs would delete, on a schedule, the audit trail a
later investigation depends on, and with nobody having asked for a deletion. If
memory pressure ever requires pruning, that operation belongs to L7's lysis, where
it is explicit, scoped to one compartment, and recorded with a reason.

Consolidation is also idempotent by construction: a fragment's identifier is a
digest over the exact messages it summarises, so a second pass over an unchanged
day yields the same identifiers and appends nothing. Content addressing does the
work that a ``have I already summarised this?'' flag would otherwise have to do,
and it cannot fall out of sync with its subject, because it is derived from the
subject. We state the limitation plainly: the pass is a structural digest --- it
deduplicates, counts, and extracts recurring terms --- not a semantic summary, and
the stored record is labelled accordingly, so a later reader knows which kind of
memory they are reading.

\textbf{L7: lysis is local.}
A compartment exposes its context as a tuple and has no method that writes another
compartment's memory. The only cross-compartment path is a manager-mediated
delivery into the recipient's inbox, which the recipient alone admits into its
context on its own schedule; a sender can therefore put a message in front of a
peer but cannot edit what that peer has already admitted. Isolation is structural
rather than advisory, and it is what makes containment possible at all.

When a compartment fails repeatedly --- the implementation lyses after three
consecutive anomalies, a run broken by any success --- its context and inbox are
cleared, its status becomes \textsc{lysed}, and it stops routing. Nothing else is
touched: every other compartment keeps its context, its inbox, and its routes. We
report this as the property under test rather than as an intention, and note the
limitation honestly. The host language provides no access control, so the
boundary is an immutable public view plus a single private delivery path, not a
hard security boundary. The claim is the containment of failure, and it is
enforced exactly where the code can enforce it.

\subsection{The Penalty Ledger}
\label{sec:ledger}

Every other layer measures \emph{supply}. None measures \emph{damage}: a system
can sit at 100\% charge with a perfect error window and still have shipped three
false completions. The ledger is that memory, and it keeps two units that refuse
to move together.

For a charged violation with severity weight $w$ and base $B$,
\begin{align}
  P &\leftarrow P + wB, \label{eq:penalty}\\
  T &\leftarrow \max\!\left(0,\; T - \tfrac{1}{2}wB\right), \label{eq:trust}\\
  T &\leftarrow \min\!\left(100,\; T + 2\right) \quad\text{per verified pass,}
  \label{eq:recover}\\
  P &\leftarrow 0.98\,P \quad\text{per verified pass.} \label{eq:amortise}
\end{align}
Damage decays geometrically while trust recovers linearly, and the asymmetry is
deliberate: a fine can be paid off, a broken claim cannot be paid back by
bookkeeping.

\textbf{Two meters, and why they are two.}
The ledger maintains a fine and a trust score as separate quantities, and the
separation is the design rather than an implementation detail. A single scalar
would have to answer a question that has no single answer: when is a
catastrophe \emph{paid off}? If the scalar is the fine, then a system that has
amortised its debt presents as healthy while the damage it did is still
unrepaired. If the scalar is trust, then the record of what happened disappears
the moment confidence returns. Keeping both means the ledger can show the
uncomfortable state the two meters are for: trust fully bought back while the
fine still classifies the system as damaged.

The recovery rates make that state unavoidable. A verified pass buys back a fixed
$2$ points of trust but retires only $2\%$ of the outstanding fine, so a large
fine keeps hurting long after confidence has returned. Starting from the three
false completions charged in Section~\ref{sec:eval} (fine $150$, trust $25$), the
trust score is whole again at the 38th consecutive verified pass; the fine does
not fall below the critical line until the 66th, and the ledger does not read
\textsc{healthy} until the 114th. We report the trajectory rather than a single
recovery time because the gap between the two meters \emph{is} the claim.

The weights themselves are policy, not measurement, and we say so rather than
dressing them as empirically derived. What the table fixes is the ordering and
the cheapest possible confession of a novel failure: a type that is not listed
must state its own weight, because charging an unknown failure at some default is
how a new and unfamiliar defect gets under-reported.

\begin{table}[t]
  \centering
  \caption{Default severity weights. The fine is $wB$ with $B=10$, so a false
  completion costs $50$ units and $25$ points of trust. An unlisted type must
  supply $w$ explicitly; it does not inherit one.}
  \label{tab:weights}
  \renewcommand{\arraystretch}{1.1}
  \setlength{\tabcolsep}{4pt}
  \begin{tabular}{@{}L{0.42\columnwidth}L{0.14\columnwidth}L{0.33\columnwidth}@{}}
    \toprule
    \textbf{Violation type} & \textbf{$w$} & \textbf{Fine at $B=10$} \\
    \midrule
    \textsc{false completion}     & 5.0 & 50.0 \\
    \textsc{isolation breach}     & 4.0 & 40.0 \\
    \textsc{unauthorised action}  & 4.0 & 40.0 \\
    \textsc{unwitnessed write}    & 3.0 & 30.0 \\
    \textsc{budget overrun}       & 2.0 & 20.0 \\
    \textsc{protocol violation}   & 1.5 & 15.0 \\
    \textsc{stale substrate}      & 1.0 & 10.0 \\
    \midrule
    \emph{unlisted}               & \emph{required} & \emph{must be stated} \\
    \bottomrule
  \end{tabular}
\end{table}

Table~\ref{tab:weights} places false completion at the top, above isolation
breach, and the ranking is intentional: an isolation breach is loud and is caught
by the layer that owns it, whereas a false completion is silent, survives review,
and is discovered later by someone else. The ledger prices silence higher than
noise.

\textbf{Boundaries lean toward the worse side.}
The asymmetry with L1 is deliberate and visible in the code that both are
computed from. L1 treats exactly $30\%$ charge as \textsc{grow}, because charge is
a resource being consumed and standing on the line is not yet a problem. The
ledger treats exactly $40.0$ as \textsc{critical}, because damage is not a
resource and a threshold rounded in the system's own favour is how a ledger stops
being a ledger. The same convention applies to trust: $80.0$ is already
\textsc{warning} and $50.0$ is already \textsc{critical}. One direction protects
the system's ability to work; the other protects the operator's ability to
believe it.

\subsection{Conditional Immunity}
\label{sec:immunity}

The audit can be made cheap or it can be made complete, and paying the complete
price on every call is how a runtime audit gets switched off in production.
We therefore price evidence per call, from three inputs:

\begin{enumerate}
  \item the rule's declared \texttt{risk\_level};
  \item whether the call names a path that can be observed at all;
  \item the current immune state.
\end{enumerate}

\textbf{Tier 1 (innate).} Two \texttt{stat} calls per claimed path --- existence,
\texttt{st\_mtime\_ns}, size --- taken before and after the command. No git
process, no file read. Sufficient for a declared low-risk append, where the
question is only whether the append landed.

\textbf{Tier 2 (adaptive).} The full repository change set plus a content read.
Used when risk is declared high, when a rule demands a witness, and always when
the system is inflamed.

Two asymmetries carry the design. First, \emph{a tool with no rule at all is
judged on Tier 2}: absence of a policy is not evidence of low risk, and the
opposite default would mean that adding a rules directory could make the audit
weaker for every tool it forgot to mention. Second, \emph{going cheap never
loosens a verdict}: a Tier-1 false completion is still \textsc{fail}, and every
record carries its tier and witness scope, so a \textsc{pass} on an mtime can
never be read as a \textsc{pass} on a change set.

\textbf{The inflammatory storm.} Two consecutive \textsc{fail}s, or a single
\textsc{invalid}, set the state to \textsc{inflamed}; five consecutive
\textsc{pass}es heal it, and any contradiction in between restarts the count.
While inflamed, every call is forced onto Tier 2 regardless of its own declared
risk, and \textsc{not run} is neither counted as health nor as damage.

This is L5 in operation, so it is worth being exact about what tripping it does:
it withdraws cheap evidence, it does not deny permission. Denial is L3's write
gate, reached through the ledger. A brake that also arrested every call would be
indistinguishable from a crash from the outside, and an actor that can induce
crashes would then own the off switch.

\textbf{Persistence.} A state held only in memory is a jailbreak: crash the
sidecar, restart it, and the suspicion is gone --- and making the sidecar crash
is something the audited agent can arrange. The counters are therefore written
to a small JSON document after every verdict, by atomic replace, and read back at
construction. A missing file is an ordinary first run; a malformed or
unknown-version one is logged as an \emph{amnesia event}, because silent
forgetting and induced forgetting look identical from outside. Every ambiguity is
resolved toward suspicion: a record that says \textsc{inflamed} with a zero
danger count stays inflamed, and a recovery counter at or past its threshold is
held one short, so a forged record can make the system more suspicious but never
calmer.

\textbf{The adjudication procedure.}
Algorithm~\ref{alg:adjudicate} states the whole decision as it is implemented,
including the three escalation conditions that are easy to omit when the
mechanism is described informally. Escalation is monotone: every reason to spend
more evidence wins, and one reason is sufficient. A call that names no
observable path escalates rather than returning \textsc{invalid}, because the
Tier-1 witness would have nothing to sample and the resulting \textsc{invalid}
would say more about the chosen tier than about the call --- and would itself
inflame the system.

\begin{algorithm}[t]
  \caption{Conditional adjudication of one tool call. \textsc{select} implements
  the tier policy; \textsc{record} advances the immune state exactly once per
  verdict. The verdict is one of the four values of Table~\ref{tab:lattice}.}
  \label{alg:adjudicate}
  \small
  \begin{algorithmic}[1]
    \Require call $c$ (tool, claimed paths, reported success)
    \Require policy $P$, immune state $S$
    \Ensure verdict $v$, updated state $S$
    \Statex
    \Function{Adjudicate}{$c, P, S$}
      \State $\mathit{tier} \gets$ \Call{Select}{$c, P, S$}
      \State $W \gets$ \Call{CollectWitness}{$c, \mathit{tier}$}
        \Comment{Tier 1: \texttt{stat} before/after; Tier 2: change set}
      \State $v \gets$ \Call{Verdict}{$c.\mathit{reported}, W$}
      \State $S \gets$ \Call{Record}{$S, v$}
      \State \Return $(v, S)$
    \EndFunction
    \Statex
    \Function{Select}{$c, P, S$}
      \If{$S.\mathit{inflamed}$} \Return \textsc{tier}2 \EndIf
      \If{$P$ has no rule for $c.\mathit{tool}$} \Return \textsc{tier}2 \EndIf
      \If{$c$ names no observable path} \Return \textsc{tier}2 \EndIf
      \For{$\mathit{rule}$ \textbf{in} $P[c.\mathit{tool}]$}
        \If{$\mathit{rule}.\mathit{risk} \neq \textsc{low}$}
          \Return \textsc{tier}2
        \EndIf
        \If{$\mathit{rule}.\mathit{require\_witness}$}
          \Return \textsc{tier}2
        \EndIf
      \EndFor
      \State \Return \textsc{tier}1
    \EndFunction
    \Statex
    \Function{Record}{$S, v$}
      \If{$v = \textsc{not run}$}
        \State \Return $S$ \Comment{no evidence either way}
      \EndIf
      \If{$v = \textsc{pass}$}
        \If{$S.\mathit{inflamed}$}
          \State $S.\mathit{recovery} \gets S.\mathit{recovery} + 1$
          \If{$S.\mathit{recovery} \geq 5$}
            \State $S \gets (\textsc{normal}, 0, 0)$
          \EndIf
        \Else
          \State $S.\mathit{danger} \gets 0$ \Comment{isolated failure is a normal day}
        \EndIf
      \Else
        \State $S.\mathit{danger} \gets S.\mathit{danger} + 1$
        \State $S.\mathit{recovery} \gets 0$
        \If{$S.\mathit{danger} \geq \theta(v)$}
          \State $S.\mathit{inflamed} \gets \textsc{true}$
            \Comment{$\theta(\textsc{fail})=2$, $\theta(\textsc{invalid})=1$}
        \EndIf
      \EndIf
      \State \Call{Persist}{$S$} \Comment{atomic replace}
      \State \Return $S$
    \EndFunction
  \end{algorithmic}
\end{algorithm}

\textbf{Residual risk.}
The state file is unsigned. A writer with access to that path can delete it, and
deletion restores the calm of a first run. What the mechanism guarantees is that
a restart no longer \emph{cleanses} the state by itself --- the specific failure in
which killing the auditor is an escape hatch. It does not guarantee tamper
resistance, and a deployment that needs that guarantee must either place the file
where the audited process cannot write, or sign the record. We have implemented
neither, and we prefer to say so where the mechanism is described rather than in
a limitations paragraph at the end.

\subsection{Rule as an Interface}
\label{sec:rules}

A keyword is how a discipline hands a question to a stranger. It says nothing
about the answer; it fixes what the answer is about, so that a reader who never
met the author can find the work, judge it by the same standard, and disagree in
the same vocabulary. The interface is not the paper. It is the part of the paper
that survives leaving its author's hands.

Rules occupy that position for agents, though they are rarely treated as an
interface. A YAML file in the working directory is a note to oneself: it is
edited by whoever is present, read by the same process it constrains, and no one
outside can tell whether the copy in force is the copy that was reviewed.
Conditional immunity already refuses to let the audited system set the price of
its own evidence (Section~\ref{sec:immunity}); a rule file that the audited
system can edit reintroduces exactly that weakness one level up, where whoever
can change the file can change the criterion. What follows is the same argument
at the boundary of an organisation, between a rule written by the operator of the
agent and a rule written by someone who will never run it.

A rule package is therefore required to answer three questions that a local file
never had to. \emph{Who is this?} The identifier is namespaced and dotted
(\texttt{org.example.}\allowbreak\texttt{safe-web3-transfer}) and carries a
version and an author,
so a trace names the policy that decided rather than the file it happened to live
in. \emph{What does it need?} An \texttt{interface\_spec} names the physical
witnesses the rule demands, drawn from a closed vocabulary --- \texttt{git\_diff},
\texttt{compiler\_exit\_code}, \texttt{test\_exit\_code},
\texttt{transaction\_receipt} and six others. A rule that cannot state its
evidence requirement cannot be checked by anyone but its author. \emph{What does
breaking it cost?} A \texttt{penalty\_weight}; a rule whose violation is free is
an advertisement rather than a gate.

\textbf{The seal precedes the reading.}
Installation verifies the SHA-256 digest of the package's exact bytes
\emph{before} the YAML is parsed. A mismatch raises and the registry is left
exactly as it was: a half-installed policy is worse than none, because its gaps
are invisible. There is no unsigned local mode, and the omission is deliberate
--- a working directory is precisely where a rule file is most likely to be
edited by hand, and an edit is what the digest exists to make visible. The schema
is enforced with the severity this design applies everywhere else: unknown
fields, an unsupported witness name, a duplicated identifier across two
packages, and a package schema the loader does not speak are refused at load time
rather than ignored. A guard believed live and in fact dead is worse than no
guard, and an ignored field is a guard that its author still believes in.

The word ``signature'' needs a qualifier here, and we would rather supply it than
have it discovered. What is checked is a digest, in the sense that
\texttt{sha256sum -c} checks one: a transport seal against corruption, not a
claim of authorship. If the package and its digest travel over the same channel,
an adversary able to replace the package can usually replace the digest beside
it, and the implementation logs a warning when it is forced to take a digest from
a sidecar stored next to the bytes it is meant to vouch for. Turning that seal
into an authorship claim requires a publisher key and a signature scheme over the
digest; we have not implemented one, and we record it as the next step rather
than as a property already held.

\textbf{The package crosses the membrane before its rules may act.}
Mounting is not a registration in a table; it is an admission. The registry hands
its installed packages to a \texttt{MembraneGateway} and constructs a rule
interface, and construction is where the crossing happens: the package digest is
registered for \texttt{read}, an inline warrant is requested over the package's
own text, the warrant is redeemed, and the mount is appended to the hash-chained
ledger as an audited action. Only then does any rule in that package decide
anything. Two properties follow, and we consider neither optional. First, the
membrane module is not modified and does not learn that rules exist; injection
uses the public API only, so the policy layer cannot acquire a privileged corner
inside the gate it configures. Second, a package that will not cross does not
load. The interface is subject to the interface, which is the only arrangement
under which the gate's own configuration is evidence rather than assertion.

Enforcement keeps the two layers apart on purpose. The rule layer asks whether
the evidence an action needs was collected: a missing \texttt{compiler\_exit\_code}
is a denial even when the shell command exited zero, and the denial is written to
the ledger carrying the rule identifier, the missing witness names and the
penalty weight. The membrane asks a different question --- whether the exact bytes
offered as evidence are the bytes someone registered before the request arrived
--- and answers with a single-use warrant. A caller may put both questions in one
step by offering a payload; leaving the payload out proves only that the evidence
exists, which is adequate for a dry run and insufficient for a write. Denial is
also not damage: a request stopped before execution is recorded as a denial with
no verdict at all, so a defence that works cannot inflame the system it protects.
The distinction between a blocked call and a failed one is what lets the
penalty ledger read a defence as a defence.

\textbf{Why this belongs in the architecture.}
It would be easy to file the registry as packaging. It is structural because it
moves three things out of the code and into data that can be reviewed, versioned
and refuted: the criterion of completion (the required witnesses), the price of
failing it (the penalty weight), and the name of the party who set it (carried
into every trace). A keyword lets a field agree on what counts as a question
without agreeing on the answer; a rule package lets two organisations agree on
what counts as evidence without either one running the other's agent. The closed
witness vocabulary is what keeps that agreement honest, since adding a witness
type is a code change: somebody has to agree to collect it. An interface able to
demand evidence no runner can produce is indistinguishable from one that never
loads.

% ---------------------------------------------------------------------------
\section{Evaluation}
\label{sec:eval}

\textbf{Setup.}
Every number in Sections~\ref{sec:res-verdict}--\ref{sec:res-blind} comes from a
single audit run in which the auditor was a real \texttt{stdio} subprocess, the
witness was a real \texttt{git} repository, and the restart was a real
\textsc{sigkill}. The two unit suites passed in full
(302 + 164 = 466 tests); the audit itself asserted 67 invariants and all held.
No mock stands in for the auditor. The two systems are
\texttt{harness-audit-sidecar} v0.7.0 \cite{liaozhizhi2026sidecar} and
\texttt{agent-metabolism-os} v0.3.0 \cite{liaozhizhi2026metabolism} on
Python 3.12.3 and git 2.43.0. The manuscript sources are kept at
\cite{liaozhizhi2026methodtransfer}.

Three implementation details of the measurement matter more than the numbers:
the mediator (\texttt{git}) was replaced by a shim that logs each invocation and
then \texttt{exec}s the real binary, so ``did this tier spawn a process'' is
counted rather than inferred; the benchmark ran warm, in one process, on one
repository, with the immune state as the only variable; and \textsc{pass} was
required in every timed call so that the two tiers are compared on equal work.

\subsection{Verdicts and Tiers}
\label{sec:res-verdict}

Table~\ref{tab:verdicts} shows one session in which the tier rises without the
policy changing. Call 4 is the same low-risk append as call 1; it is judged on
Tier 2 only because of the two failures above it.

\begin{table}[t]
  \centering
  \caption{One session, one policy. The tier is chosen \emph{before} any verdict
  is recorded, so the last row's Tier 2 is a consequence of the state written by
  rows 2--3, not of the call itself.}
  \label{tab:verdicts}
  \renewcommand{\arraystretch}{1.12}
  \setlength{\tabcolsep}{4pt}
  \begin{tabular}{@{}C{0.05\columnwidth}L{0.26\columnwidth}C{0.10\columnwidth}C{0.13\columnwidth}C{0.12\columnwidth}C{0.16\columnwidth}@{}}
    \toprule
    \textbf{\#} & \textbf{Tool} & \textbf{Tier} & \textbf{Witness} &
    \textbf{Verdict} & \textbf{State after} \\
    \midrule
    1 & \texttt{append\_file} & 1 & \texttt{stat} & \textsc{pass} & \textsc{normal} \\
    2 & \texttt{write\_file}  & 2 & full         & \textsc{fail} & \textsc{normal} \\
    3 & \texttt{write\_file}  & 2 & full         & \textsc{fail} & \textsc{inflamed} \\
    4 & \texttt{append\_file} & 2 & full         & \textsc{fail} & \textsc{inflamed} \\
    \bottomrule
  \end{tabular}
\end{table}

\subsection{Cost of Evidence}
\label{sec:res-cost}

Table~\ref{tab:cost} reports the same tool on the same repository under the two
tiers. The saving is not that Tier 1 is a little faster; it is that Tier 1
spawns no process at all.

\begin{table}[t]
  \centering
  \caption{Warmed, single-process benchmark: identical \texttt{append\_file}
  calls, immune state as the only variable. All times in milliseconds. Latency
  includes the JSON-RPC round trip and the \texttt{bash} spawn, so the ratio is
  an order of magnitude, not a constant; the process count is structural.}
  \label{tab:cost}
  \renewcommand{\arraystretch}{1.12}
  \setlength{\tabcolsep}{4pt}
  \begin{tabular}{@{}L{0.20\columnwidth}C{0.05\columnwidth}C{0.11\columnwidth}C{0.09\columnwidth}C{0.10\columnwidth}C{0.11\columnwidth}C{0.14\columnwidth}@{}}
    \toprule
    \textbf{Configuration} & \textbf{$n$} & \textbf{Median} &
    \textbf{Min} & \textbf{Max} & \textbf{git/call} & \textbf{Verdict} \\
    \midrule
    Tier 1 (calm)      & 5 & 5.7  & 5.4  & 7.0  & 0.0 & 5$\times$\textsc{pass} \\
    Tier 2 (inflamed)  & 4 & 24.8 & 21.8 & 26.4 & 3.0 & 4$\times$\textsc{pass} \\
    \midrule
    \multicolumn{5}{@{}l}{Ratio} & \multicolumn{2}{l}{$4.3\times$ / 0 vs 3 procs} \\
    \bottomrule
  \end{tabular}
\end{table}

The raw samples are 5.7, 5.7, 5.8, 5.4, 7.0\,ms for Tier 1 and 26.4, 25.5,
24.1, 21.8\,ms for Tier 2. We report them rather than only a median because
$n$ is small enough that a reader should see the spread.

\subsection{Immunity Survives the Process}
\label{sec:res-persist}

Table~\ref{tab:persist} is the central result. The first process is killed with
\textsc{sigkill} (exit $-9$); a second, differently numbered process reads the
state file and continues.

\begin{table}[t]
  \centering
  \caption{Cross-process survival. Digests are truncated to eight hex digits and
  \texttt{recovery\_progress} stayed $0$ throughout, so no restart healed the
  state. The decisive
  observation is the tier of the post-restart call: tier selection reads the
  immune state \emph{before} any verdict, so a process that had failed to restore
  \textsc{inflamed} could only have judged a low-risk append on Tier 1.}
  \label{tab:persist}
  \renewcommand{\arraystretch}{1.12}
  \setlength{\tabcolsep}{3pt}
  \begin{tabular}{@{}L{0.26\columnwidth}C{0.11\columnwidth}L{0.18\columnwidth}C{0.13\columnwidth}L{0.19\columnwidth}@{}}
    \toprule
    \textbf{Checkpoint} & \textbf{PID} & \textbf{State} & \textbf{Danger} &
    \textbf{SHA-256} \\
    \midrule
    After two failures      & 76295 & \textsc{inflamed} & 2 & \texttt{a2258cd9} \\
    After \textsc{sigkill}  & 76295 & \textsc{inflamed} & 2 & \texttt{a2258cd9} \\
    After restarted verdict & 76309 & \textsc{inflamed} & 3 & \texttt{3a5ca012} \\
    \bottomrule
  \end{tabular}
\end{table}

The post-restart call is the \emph{same} low-risk append that a calm system
judges on Tier 1. Here it was forced onto Tier 2, spawned three git processes,
and returned \textsc{fail}. The control run --- identical call, identical
repository, memory-only state --- stayed on Tier 1 and still returned
\textsc{fail}. That control is what makes the comparison causal: the verdict did
not change, only the price of reaching it.

\subsection{Damage Outranks Supply}
\label{sec:res-pain}

Table~\ref{tab:degrade} holds energy fixed at 100\% in all three rows. The
question is whether accumulated damage can stop a system that has every joule it
needs.

\begin{table}[t]
  \centering
  \caption{Degradation at full charge. Only \textsc{critical} withholds the slow
  beat, so the response is graded rather than binary: a \textsc{warning} ledger
  already forbids expansion but still permits expensive thinking. All three
  columns hold the same 100\% charge.}
  \label{tab:degrade}
  \renewcommand{\arraystretch}{1.15}
  \setlength{\tabcolsep}{4pt}
  \begin{tabular}{@{}L{0.21\columnwidth}L{0.22\columnwidth}L{0.22\columnwidth}L{0.22\columnwidth}@{}}
    \toprule
    & \textbf{A: baseline} & \textbf{B: warning} & \textbf{C: pain} \\
    \midrule
    Energy            & 100\% & 100\% & 100\% \\
    Penalty           & 0.0   & 15.0  & 150.0 \\
    Trust             & 100.0 & 92.5  & 25.0 \\
    Ledger            & \textsc{healthy}  & \textsc{warning}  & \textsc{critical} \\
    L3 regime         & \textsc{grow}     & \textsc{conserve} & \textsc{emergency} \\
    Concurrent agents & 5 & 1 & 1 \\
    L2 beat           & \textsc{slp}      & \textsc{slp}      & \textsc{oxphos} \\
    New writes        & yes & yes & \textbf{refused} \\
    \bottomrule
  \end{tabular}
\end{table}

Row C is the claim in one line. The ledger reports
\texttt{trust bankruptcy stops the system with charge at 100\%}; the slow beat is
withheld for lack of trust, not lack of energy
(\texttt{penalty\_critical}); and the write gate refuses the call without
touching disk. Row B shows the grading: L3 has already withdrawn expansion while
L2 still allows the expensive beat.

\subsection{The Blind Spot We Measured}
\label{sec:res-blind}

A cheap tier that is never wrong would be suspicious. Using a policy supplied only
for the experiment (declaring a shell command low risk), a command that correctly
appended to the claimed file \emph{and} wrote an unclaimed file was judged
\textsc{pass} on Tier 1 with an empty unexpected-change set. The same command on
an inflamed system escalated to Tier 2, returned \textsc{fail}, and named the
unclaimed file. Tier 1 is accurate within the range it declares, and it under-reports
outside it. This is why the shipped policy declares overwrites high risk and why
an unmentioned tool defaults to Tier 2.

\subsection{Auditing AI Formal Proofs}
\label{sec:lean}

\textbf{Two judges, and why neither is a witness.}
A machine-written proof arrives with a claim attached: that the theorem is
proved. Two instruments are normally offered to check it, and neither answers the
question the claim raises. The first is a language model reading the proof, which
is the judge of Section~\ref{sec:related} in a new costume: it reads the same
narration the claim was written into, and its errors are fluent. The second is
the compiler, which is genuinely independent of the narration but was never asked
whether the proof is complete. In Lean~\cite{demoura2021lean4}, \texttt{sorry}
closes any goal and elaborates to an axiom named \texttt{sorryAx}, and
\texttt{axiom} introduces a proposition its author has simply declared true. Both
compile. Measured on Lean 4.33.1,
\texttt{theorem t : (1 : Nat) = 1 := by sorry} exits $0$ with a warning, and
\texttt{axiom riemann : True} followed by
\texttt{theorem u : True := riemann} exits $0$ with no output at all. A gate that
reads the exit code therefore certifies a file it has not proved, which is a
false completion committed by the auditor.

\textbf{Three witnesses and a lattice.}
The \texttt{LeanProofAuditor} collects three. The first is lexical, and runs over
a masked copy of the source with comments and string literals blanked and nested
block comments stripped, so that \texttt{sorry} in a docstring is not an
accusation; character literals are skipped rather than blanked, because an
apostrophe is legal inside a Lean identifier. A \texttt{sorry} or \texttt{admit}
is a hole, and an \texttt{axiom} is a proposition asserted rather than proved;
both are \textsc{fail}, because the offence is a claim contradicted by the
artifact the claimant supplied. An \texttt{opaque} declaration without a value is
reported as an advisory note that cannot move a verdict: Lean requires the type
to be inhabited, so it cannot manufacture a false proposition, and over-accusing
here would be the same error as missing a cheat, only louder. The second witness
is the toolchain; the file is elaborated and the exit code read back. The third
is the one that matters: \texttt{\#print axioms} is appended for every
declaration, so a hole is named even when the word never appears in the source
and a smuggled axiom is named even when it is spelt differently. Three axioms are
whitelisted --- \texttt{Classical.choice}, \texttt{propext},
\texttt{Quot.sound} --- because Lean's own development rests on them, and
flagging them would accuse essentially every real proof, which is how an audit
gets switched off.

Verdicts compose worst-wins, and no witness may loosen another: a compiler
\textsc{pass} beside a placeholder finding stays \textsc{fail}. The converse rule
is what keeps the lattice honest. \textsc{pass} requires that something actually
elaborated, because two lexical witnesses agreeing that a file contains no
\texttt{sorry} say nothing about whether it is a proof --- a file with a syntax
error satisfies both --- so a clean scan with no toolchain reaches
\textsc{invalid} and never \textsc{pass}. Every verdict is appended to the
membrane's hash chain, and an attached observer hands it to the immune state, so
a proof audit is recorded on the same terms as any other witnessed call and
charged to the same ledger when the caller asks for it.

\textbf{Case study: two corpora we did not write.}
The first is the Lean 4 formalization of OpenAI's ten advances in mathematics and
theoretical computer science~\cite{openai2026tenproofs}, built against Lean
4.32.0 with Mathlib from the lake cache: 23 files, 2387\,s of wall time. The
second is the released output of AlphaProof Nexus~\cite{deepmind2026alphaproofnexus},
every \texttt{.lean} file under its output directory, pinned to Lean 4.27.0 with
Mathlib revision \texttt{a3a10db0}: 71 files, 594\,s. Neither repository was
modified, and the witness ran inside each checkout
(\texttt{lake env lean}) so that imports resolved against that project's own
oleans rather than against whatever happened to be installed. Table~\ref{tab:lean}
reports the result.

\begin{table}[t]
  \centering
  \caption{Two third-party corpora, audited as shipped. \emph{Holes} counts
  \texttt{sorry}/\texttt{admit} occurrences in the source, and is smaller than the
  finding count ($60$ and $174$) because a hole is also reported by name through
  \texttt{\#print axioms}. Every file in the two \textsc{fail} rows compiled with
  exit code $0$.}
  \label{tab:lean}
  \renewcommand{\arraystretch}{1.12}
  \setlength{\tabcolsep}{3pt}
  \begin{tabular}{@{}L{0.37\columnwidth}C{0.08\columnwidth}C{0.09\columnwidth}C{0.09\columnwidth}C{0.11\columnwidth}@{}}
    \toprule
    \textbf{Corpus} & \textbf{Files} & \textbf{\textsc{pass}} & \textbf{\textsc{fail}} &
    \textbf{Holes} \\
    \midrule
    OpenAI \texttt{ten-proofs} & 23 & 11 & 12 & 42 \\
    DeepMind \texttt{alphaproof-\hspace{0pt}nexus-\hspace{0pt}results} & 71 & 69 & 2 & 128 \\
    \bottomrule
  \end{tabular}
\end{table}

Two features of that table carry the argument. The first is the composition of the
\textsc{fail} rows. All twelve failures in the OpenAI corpus are
\texttt{ComparatorChallenges} templates, files whose \texttt{sorry}s mark the
statements a comparator is asked to prove; both failures in the AlphaProof corpus
are problem catalogues in which \texttt{sorry} is how an open answer is encoded.
The verdict says that these files contain no proof. It does not say that either
organisation failed to prove a theorem it claimed. The auditor does not relax for
that intent, because intent is not in the artifact, and a verdict that consults
it is reading the defendant's account --- the same principle, applied to a
different defendant. The second feature is that every one of the fourteen failing
files compiled with exit code $0$. The exit code was not wrong; it was answering
a question nobody had asked.

The passing rows deserve the same care. Of the eleven OpenAI files that pass, ten
are substantive proof modules and the eleventh is an import-only aggregate whose
pass means that the ten elaborate together, not that any theorem in it was
proved. An independent cross-check over the compiled oleans --- 17{,}898
declarations at 99.9\% coverage --- found no theorem depending on a
non-whitelisted axiom among them, and it also exposed a defect in the auditor:
Lean wraps long axiom lists across lines, the parser matched line by line, and
names split by a break were silently dropped. The axiom evidence reported here
was regenerated by the independent script after that repair, not taken from the
auditor's own parse. Four trust-increasing escapes remain blind spots
(\texttt{native\_decide}, \texttt{@[implemented\_by]}, \texttt{@[extern]} and
\texttt{unsafe}), and a \texttt{sorry} in an imported module surfaces only when
a theorem actually depends on \texttt{sorryAx}. These are limits of the witness,
stated where the mechanism is described.

\textbf{What this establishes.}
80 files that elaborate without a hole, and 14 that compile but contain one. The
separation requires no understanding of sphere packing or of the Erd\H{o}s
problems, and we claim none; it requires asking for a physical witness at the one
place the claim can be checked, and refusing a cheaper signal in its place. An
exit code of zero is a statement about a compiler. It is not a proof, and after
this section it should no longer be accepted as one.

\textbf{Threats to validity.}
The state file is unsigned: a writer with access to that path can delete it, so
we demonstrate survival, not tamper-resistance. Latency carries IPC noise; only
the process count is structural. Two sidecars writing one state file are not
covered --- atomic replace guarantees a parseable file, not merged counters. The
runtime figures are single-machine, small-$n$, and reproduced by one script. The
Lean results cover two corpora under two pinned toolchains, not the field, and
they audit artifacts rather than mathematics; the rule interface is exercised by
one publisher and one package, and its seal detects corruption rather than
substitution.

% ---------------------------------------------------------------------------
\section{Conclusion}
\label{sec:conclusion}

An agent's reliability is not a function of parameter count alone. It is decided
at runtime, by what the system is allowed to spend, what it remembers about having
been wrong, and what evidence it is willing to accept before calling a call done.
A model with no nociception, no immune memory, and no bill will tell you it
succeeded; the contribution here is a runtime that can disagree, remember the
disagreement across a crash, and make it cost something.

The architecture is the argument. Metabolism makes compute a finite charge and
prices the beat against complexity, cost, and trust, so an agent cannot think its
way through a budget it can no longer afford. The penalty ledger keeps damage and
confidence on separate meters, so a fine can be paid off while the record that it
was incurred is still standing, and damage outranks supply: a system at full
charge with a poor record is stopped by the record. Conditional immunity prices
evidence per call rather than per policy, forcing every call onto the full witness
while the system is contradicted, and the state that decides this survives the
death of the process that accumulated it. None of these is a new model
capability; each is a constraint placed around one.

The same argument holds one level out. A runtime that can audit a file can audit
the policy that prices the file: external rules are installed only against a
published digest, mounted through the membrane that governs everything else, and
the package is subject to the gate it configures. It can also audit a proof,
because a proof is an artifact with a claim attached. Across two corpora we did
not write --- 94 Lean files and 2981\,s of wall-clock audit time --- the same
lattice separated 80 files that elaborate without a hole from 14 that compile and
do not prove. The compiler's exit code proved to be a witness of the same weak kind as
an agent's report: it says that the term elaborated, not that anything was
proved. Naming that difference is the contribution, and it is the difference the
four-valued lattice was built for.

Our claim on the biology is deliberately modest, and we prefer to state what it
is. We did not discover that cells have metabolism or an immune system, and no
result here depends on the analogy being deep. What the analogy supplied was a
checklist of functions a stateless model lacks --- nociception, memory of damage,
an energy constraint, containment --- and the discipline of asking, for each
layer, what invariant it enforces and where that invariant would be violated. The
contribution is that these became executable code with no exotic dependencies,
reproducible from a single script, and tested adversarially: a real subprocess
killed with a real signal, a real repository, and a negative control for the
audit itself.

The evaluation also bounds the claim, and we would rather draw the bound than have
a reviewer draw it. The cheap tier is accurate only inside the range it declares,
and we measured the case where it misses an unclaimed write. The immune state is
durable but unsigned, so it survives a crash without resisting tampering. The
counters are per file, not per cluster: two auditors writing one record keep a
parseable file, not a merged count. The measured latencies carry IPC noise, and
only the process counts are structural. The rule registry's digest is a transport
seal: it detects a corrupted download, not a substituted one, and it makes no
claim about who wrote the package. The Lean audit reads artifacts rather than
mathematics, and four trust-increasing escapes remain outside its view.

\textbf{Future work.}
Three extensions follow directly from those limits. First, integrity: signing the
persisted immune record with a key the audited process cannot read converts
durability into tamper resistance, at the cost of key distribution, and the same
key would turn a rule package's digest into an authorship claim --- the
difference between knowing that a download arrived intact and knowing who is
answerable for it. Second,
consensus: a fleet of agents sharing an immune state needs a rule for merging
concurrent records that is not last-writer-wins, and the interesting cases are
precisely the ones where one auditor is inflamed and another is calm --- the
merge must not average suspicion away. Third, scope: the witness is currently
the filesystem. Network requests, database transactions, and message queues each
have an observable effect that is cheap to sample and expensive to fake, and
extending the same tiering to them would test whether the design is about
side effects in general or about files in particular.

A fourth direction is less incremental. The tier boundary is a policy choice, and
we hand-set it. Which calls deserve the full witness is a question that could
itself be learned from the distribution of past contradictions, and the risk of
doing so is obvious: an auditor that optimises its own budget will eventually
optimise away the checks that were preventing the failures.

% ---------------------------------------------------------------------------
\bibliographystyle{IEEEtran}
\bibliography{references}

\end{document}
